% 复合图: 图1标识墩立体示意(正四棱锥P-EFGH + 长方体ABCD-EFGH, 虚线为隐藏线, 含原图侧视/正视观察方向箭头)
% + 图2正视图(下矩40x20, 上三角形底40高60) + 图3俯视图(40x40正方形带对角线).
% 缺失的侧(左)视图由考生作答, 不补造; 立体透视仅示意; 横向并排, 图名基线对齐.
\begin{tikzpicture}[line width=.6pt,font=\normalsize,x=1cm,y=1cm,>=Stealth]
% 图1 立体(示意, 藏线用虚线)
\begin{scope}[shift={(0,1.2)}]
% 长方体前视面
\draw (0,0) -- (2,0) -- (2,1) -- (0,1) -- cycle;
% 深度方向
\draw[dashed] (0,0) -- (0.5,0.5);
\draw (2,0) -- (2.5,0.5);
\draw[dashed] (0.5,0.5) -- (2.5,0.5);
\draw[dashed] (0.5,0.5) -- (0.5,1.5);
\draw (2.5,0.5) -- (2.5,1.5);
\draw (0,1) -- (2,1);
\draw[dashed] (0,1) -- (0.5,1.5);
\draw (2,1) -- (2.5,1.5);
\draw[dashed] (0.5,1.5) -- (2.5,1.5);
% 棱锥
\draw (0,1) -- (1.25,4.25);
\draw (2,1) -- (1.25,4.25);
\draw (2.5,1.5) -- (1.25,4.25);
\draw[dashed] (0.5,1.5) -- (1.25,4.25);
% 顶点标签
\node[above] at (1.25,4.25) {$P$};
\node[left] at (0,1) {$E$};
\node[right] at (2,1) {$F$};
\node[right] at (2.5,1.5) {$G$};
\node[left] at (0.5,1.5) {$H$};
\node[below left] at (0,0) {$A$};
\node[below right] at (2,0) {$B$};
\node[right] at (2.5,0.5) {$C$};
\node[left] at (0.5,0.5) {$D$};
% 观察方向示意(原图所有)
\draw[->] (-1.2,1) -- (-0.5,1);
\node[left] at (-1.2,1.3) {侧视};
\draw[->] (-.2,-.8) -- (.6,-.15);
\node[below] at (-.5,-1.05) {正视};
\end{scope}
\node at (1.25,-0.7) {\bitmapinsertcaption[1]};
% 图2 正视图: 矩40x20 + 三角底40高60 (1cm=20cm)
\begin{scope}[shift={(4.2,0.6)}]
\draw (0,0) rectangle (2,1);
\draw (0,1) -- (2,1) -- (1,4) -- cycle;
% 底部尺寸 40
\draw[thin] (0,0) -- (0,-0.5);
\draw[thin] (2,0) -- (2,-0.5);
\draw[<->] (0,-0.3) -- (2,-0.3);
\node[below] at (1,-0.3) {40};
% 右侧矩高 20
\draw[thin] (2,0) -- (2.5,0);
\draw[thin] (2,1) -- (2.5,1);
\draw[<->] (2.3,0) -- (2.3,1);
\node[right] at (2.3,0.5) {20};
% 右侧锥高 60
\draw[thin] (1,4) -- (2.5,4);
\draw[thin] (2,1) -- (2.5,1);
\draw[<->] (2.3,1) -- (2.3,4);
\node[right] at (2.3,2.5) {60};
\end{scope}
\node at (5.2,-0.7) {\bitmapinsertcaption[2] 正（主）视图};
% 图3 俯视图: 40x40正方形 + 对角线
\begin{scope}[shift={(8.0,0.6)}]
\draw (0,0) rectangle (2,2);
\draw (0,0) -- (2,2);
\draw (0,2) -- (2,0);
% 底部尺寸 40
\draw[thin] (0,0) -- (0,-0.5);
\draw[thin] (2,0) -- (2,-0.5);
\draw[<->] (0,-0.3) -- (2,-0.3);
\node[below] at (1,-0.3) {40};
% 右侧尺寸 40
\draw[thin] (2,0) -- (2.5,0);
\draw[thin] (2,2) -- (2.5,2);
\draw[<->] (2.3,0) -- (2.3,2);
\node[right] at (2.3,1) {40};
\end{scope}
\node at (9.0,-0.7) {\bitmapinsertcaption[3] 俯视图};
\end{tikzpicture}
